\section{第四讲：Hitchin 纤维化}

20 世纪 80 年代后半期，Hitchin 在研究规范场论时引入了著名的可积系统，即 Higgs 丛的模空间。差不多同时，Kazhdan 与 Lusztig 引入仿射 Springer 纤维，作为 Springer 纤维的自然类比。此后十余年里这两个对象一直彼此无关，直到 B-C. Ng\^o 看出两者之间的联系。Ng\^o 的基本洞见可以概括为：Hitchin 纤维是仿射 Springer 纤维的整体类比，而仿射 Springer 纤维是 Hitchin 纤维的局部版本。这里的「整体」指涉及整体函数域或代数曲线的对象，而不只是局部函数域或形式圆盘。这一观察连同精巧的技术工作，使 Ng\^o 证明了 Langlands 与 Shelstad 所猜想的轨道积分的基本引理。下面在略为一般的设定中回顾 Hitchin 的可积系统，并精确说明它与仿射 Springer 纤维的联系。

\subsection{Hitchin 模叠}

\subsubsection{设定} 回到 \S\ref{ss:notation} 的设定。此外，固定 $k$ 上一条代数曲线 $X$（假定 $k$ 代数闭），它光滑、射影且连通。

\subsubsection{向量丛的模叠}
存在模叠 $\Bun_{n}$，分类 $X$ 上秩 $n$ 的向量丛。对任意 $k$-代数 $R$，$\Bun_{n}(R)$ 是 $X_{R}:=X\times_{\Spec k}\Spec R$ 上秩 $n$ 向量丛（局部自由凝聚层）的群胚。叠 $\Bun_{n}$ 是代数叠，见 \cite[定理 4.6.2.1]{LMB}。此外，它在 $k$ 上光滑且局部有限型。

\subsubsection{$G$-主齐性空间}
回忆 $X$ 上的{\em（右）$G$-主齐性空间}是指 $X$ 上概形 $\calE$，带有逐纤维的 $G$-作用，使得在 $X$ 的 étale 拓扑下局部地，$\calE$ 变成 $G\times X$，而 $G$-作用变成 $G$ 在第一个因子上的右平移作用。

对一般约化群 $G$，有 $X$ 上 $G$-主齐性空间的模叠 $\Bun_{G}$。对 $k$-代数 $R$，$\Bun_{G}$ 的 $R$-点是 $X_{R}$ 上 $G$-主齐性空间的群胚。于是 $\Bun_{G}$ 也是 $k$ 上光滑且局部有限型的代数叠。

\subsubsection{相伴丛}\label{sss:ass bun} 设 $(V,\rho)$ 是 $G$ 的 $k$-表示，$\calE$ 是 $X$ 上 $G$-主齐性空间。则存在 $X$ 上向量丛 $\rho(\calE)$，其全空间为
\begin{equation*}
\textup{Tot}(\rho(\calE))=G\bs (\calE\times V),
\end{equation*}
其中 $G$ 通过 $g\cdot(e,v)=(eg^{-1}, \rho(g)v)$ 作用在 $(e,v)\in\calE\times V$ 上。称向量丛 $\rho(\calE)$ 相伴于 $\calE$ 与 $\rho$。

当 $G=\GL_{n}$ 时，有群胚的等价
\begin{equation}\label{GLn vec}
\{ \mbox{$X$ 上秩 $n$ 的向量丛 $\calV$}\}\cong\{\mbox{$X$ 上的 $\GL_{n}$-主齐性空间 $\calE$}\}.
\end{equation}
方向 $\calV\mapsto\calE$ 把向量丛 $\calV$ 送到 $\calV$ 的标架构成的 $\GL_{n}$-主齐性空间，即取 $\calE=\un{\Isom}_{X}(\calO^{n}_{X},\calV)$，附带 $\GL_{n}$ 在平凡丛 $\calO^{n}_{X}$ 上的自然作用。另一方向 $\calE\mapsto\calV$ 把 $X$ 上 $\GL_{n}$-主齐性空间 $\calE$ 送到相伴于 $\calE$ 与 $\GL_{n}$ 的标准表示 $\St$ 的向量丛 $\St(\calE)$。等价 \eqref{GLn vec} 给出叠的典范同构 $\Bun_{n}\cong\Bun_{\GL_{n}}$。

\subsubsection{}\label{sss:Bun Sp} 对其他典型群 $G$，$G$-主齐性空间有更显明的向量丛描述。例如当 $G=\SL_{n}$ 时，$X$ 上 $G$-主齐性空间等同于数对 $(\calV,\iota)$，其中 $\calV$ 是 $X$ 上秩 $n$ 的向量丛，$\iota$ 是线丛同构 $\wedge^{n}\calV\cong\calO_{X}$。

当 $G=\Sp_{2n}$ 时，$X$ 上 $G$-主齐性空间 $\calE$ 的群胚等价于数对 $(\calV,\omega)$ 的群胚，其中 $\calV$ 是 $X$ 上秩 $2n$ 的向量丛，$\omega:\wedge^{2}\calV\to\calO_{X}$ 是凝聚层的 $\calO_{X}$-线性映射，且在每个几何纤维上给出辛形式。一个方向把 $\Sp_{2n}$-主齐性空间 $\calE$ 送到 $(\calV,\omega)$：其中 $\calV=\St(\calE)$ 是相伴于 $\calE$ 与 $\Sp_{2n}$ 的标准表示 $\St$ 的向量丛，而 $\calV$ 上的辛形式 $\omega$ 来自 $\Sp_{2n}$-表示之间的典范映射 $\wedge^{2}(\St)\to \one$（其中 $\one$ 是平凡表示）。

\subsubsection{Higgs 丛} 固定 $X$ 上线丛 $\calL$。$X$ 上{\em 秩 $n$ 的 $\calL$-扭 Higgs 丛}是指数对 $(\calV,\varphi)$，其中 $\calV$ 是 $X$ 上秩 $n$ 的向量丛，$\varphi:\calV\to\calV\otimes\calL$ 是 $\calO_{X}$-线性映射。自同态 $\varphi$ 称为 $\calV$ 上的{\em Higgs 场}。

存在模叠 $\calM_{n,\calL}$，分类 $X$ 上秩 $n$ 的 $\calL$-扭 Higgs 丛。忘掉 Higgs 场的态射 $\calM_{n,\calL}\to \Bun_{n}$ 可表。因此 $\calM_{n,\calL}$ 也是 $k$ 上代数叠。

\subsubsection{$G$-Higgs 丛} {\em $\calL$-扭 $G$-Higgs 丛}是指数对 $(\calE,\varphi)$，其中 $\calE$ 是 $X$ 上 $G$-主齐性空间，$\varphi$ 是 $X$ 上向量丛 $\Ad(\calE)\otimes\calL$ 的整体截面。这里 $\Ad(\calE)$ 是相伴于 $\calE$ 与 $G$ 的伴随表示 $(\frg,\Ad)$ 的向量丛，见 \S\ref{sss:ass bun}。称 $\varphi$ 是 $\calE$ 上的{\em $\calL$-扭 Higgs 场}。

当 $G=\GL_{n}$ 时，$\calL$-扭 $G$-Higgs 丛的概念等价于秩 $n$ 的 $\calL$-扭 Higgs 丛。事实上，对每个 $\calL$-扭 $G$-Higgs 丛 $(\calE,\varphi)$，得到 Higgs 丛 $(\calV=\St(\calE),\phi)$；其中把 $\phi:\calV\to\calV\otimes\calL$ 看作 $\underline{\End}(\calV)\otimes\calL$ 的整体截面，它在典范同构 $\underline{\End}(\calV)\cong\Ad(\calE)$ 下对应 $\varphi$。

还有 $X$ 上 $\calL$-扭 Higgs $G$-主齐性空间的模叠 $\calM_{G,\calL}$。$\calM_{G,\calL}$ 的 $R$-点是 $X_{R}$ 上 $\calL_{R}$-扭 $G$-Higgs 丛的群胚，其中 $\calL_{R}$ 表示 $\calL$ 到 $X_{R}$ 的拉回。遗忘态射 $\calM_{G,\calL}\to \Bun_{G}$ 可表，故 $\calM_{G,\calL}$ 是 $k$ 上代数叠。当 $G=\GL_{n}$ 时，有叠的典范同构 $\calM_{n,\calL}\cong\calM_{\GL_{n},\calL}$。

\subsubsection{例子} 当 $G=\SL_{n}$ 时，$X$ 上 $\calL$-扭 $G$-Higgs 丛等同于三元组 $(\calV,\iota,\varphi)$，其中 $\calV$ 是 $X$ 上秩 $n$ 的向量丛，$\iota:\wedge^{n}\calV\isom\calO_{X}$，且 $\varphi: \calV\to\calV\otimes\calL$ 满足 $\Tr(\varphi)=0$。

当 $G=\Sp_{2n}$ 时，$X$ 上 $\calL$-扭 $G$-Higgs 丛等同于三元组 $(\calV,\omega,\varphi)$，其中 $\calV$ 是 $X$ 上秩 $2n$ 的向量丛，$\omega:\wedge^{2}\calV\to\calO_{X}$ 逐纤维非退化，且 $\varphi: \calV\to\calV\otimes\calL$ 使得对 $\calV$ 的一切局部截面 $u$ 与 $v$，$\omega(\varphi(u),v)+\omega(u,\varphi(v))=0$ 作为 $\calL$ 的局部截面成立。

\subsubsection{Hitchin 模叠作为余切丛}
Hitchin 的原始论文 \cite{Hi} 考虑了 $\calL=\omega_{X}$ 为 $X$ 上 $1$-形式层的特例。此特例尤其重要，因为 $\calM_{G,\omega_{X}}$ 与 $\Bun_{G}$ 的余切丛密切相关。对点 $\calE\in\Bun_{G}(R)$，即 $X_{R}$ 上 $G$-主齐性空间，$\Bun_{G}$ 在 $\calE$ 处的余切复形由复形 $\Ad(\calE)^{\vee}\otimes\omega_{X}$ 在 $X_{R}$ 上的导出整体截面给出。用 $\frg$ 上的 Killing 形式可以把 $\Ad(\calE)^{\vee}$ 与 $\Ad(\calE)$ 等同，因此 $\Bun_{G}$ 在 $\calE$ 处的余切复形是 $\bR\Gamma(X_{R}, \Ad(\calE)\otimes\omega_{X})$，它集中在 $0$ 与 $1$ 次。特别地，当 $\calE\in \Bun_{G}(k)$ 的自同构群有限（例如 $\calE$ 稳定）时，$\calE$ 处的{\em Zariski 余切空间}是 $\cohog{0}{X,\Ad(\calE)\otimes\omega_{X}}$；也就是说，$\Bun_{G}$ 在 $\calE$ 处的余切向量等同于 $\calE$ 上 $\omega_{X}$-扭 Higgs 场。因此 $T^{*}\Bun_{G}$（适当定义）与 $\calM_{G,\omega_{X}}$ 共享开子叠 $T^{*}\Bun^{s}_{G}$，其中 $\Bun^{s}_{G}$ 是稳定 $G$-丛的开子叠。

\subsection{Hitchin 纤维化}

\subsubsection{$\GL_{n}$ 的 Hitchin 纤维化}
对 $X$ 上 $\calL$-扭 Higgs 丛 $(\calV,\varphi)$，在 $X$ 上局部地可把 $\varphi$ 看作元素为 $\calL$ 的局部截面的矩阵，从而可取此矩阵的特征多项式。此多项式各系数是 $\calL^{\otimes i}$ 的良定义整体截面，$1\leq i\leq n$。更内蕴地说，$\varphi$ 诱导映射 $\wedge^{i}\varphi:\wedge^{i}\calV\to\wedge^{i}\calV\otimes\calL^{\otimes i}$，可取
\begin{equation*}
a_{i}(\varphi):=\Tr(\wedge^{i}\varphi)\in\cohog{0}{X,\calL^{\otimes i}}.
\end{equation*}
这样就定义了态射
\begin{equation*}
f: \calM_{n,\calL}\to\calA_{n,\calL}:=\prod_{i=1}^{n}\cohog{0}{X,\calL^{i}},
\end{equation*}
它把 $(\calV,\varphi)$ 送到 $(a_{1}(\varphi),\cdots, a_{n}(\varphi))$。把 $\calA_{n,\calL}$ 看作 $k$ 上仿射空间。态射 $f$ 称为 $G=\GL_{n}$ 情形下的{\em Hitchin 纤维化}。

\subsubsection{一般的 Hitchin 纤维化}\label{sss:Hit fib general} 对 \S\ref{ss:notation} 中一般连通约化群 $G$，$\GL_{n}$ 情形中特征多项式的系数由 $\frg$ 上 $G$-不变的基本多项式代替。回忆 $\frc=\frg\sslash G=\Spec\Sym(\frg^{*})^{G}$。Chevalley 定理说 $\Sym(\frg^{*})^{G}\cong\Sym(\frt^{*})^{W}$，而后者是 $r$ 元多项式环。固定 $\Sym(\frg^{*})^{G}$ 作为 $k$-代数的齐次生成元 $f_{1},\cdots, f_{r}$，其次数 $d_{1}\leq \cdots \leq d_{r}$ 是典范确定的，尽管 $f_{i}$ 本身不典范。当 $G$ 几乎单时，数 $e_{i}=d_{i}-1$ 是 $G$ 的{\em 指数}。把 $f_{i}$ 看作 $G$-不变对称多重线性函数 $\frg^{\otimes d_{i}}\to k$，则对 $X$ 上任一 $G$-主齐性空间 $\calE$，$f_{i}$ 诱导相伴丛之间的映射
\begin{equation*}
f_{i}: \Ad(\calE)^{\otimes d_{i}}\to\calO_{X}.
\end{equation*}
这又诱导
\begin{equation*}
f^{\calL}_{i}:(\Ad(\calE)\otimes\calL)^{\otimes d_{i}}\to\calL^{\otimes d_{i}}.
\end{equation*}
若 $\varphi$ 是 $\calE$ 上 $\calL$-扭 Higgs 场，可把 $f^{\calL}_{i}$ 作用在 $(\Ad(\calE)\otimes\calL)^{\otimes d_{i}}$ 的截面 $\varphi^{\otimes d_{i}}$ 上，得到
\begin{equation*}
a_{i}(\varphi):=f^{\calL}_{i}(\varphi^{\otimes d_{i}})\in\cohog{0}{X,\calL^{\otimes d_{i}}},  \quad i=1,2,\cdots, r.
\end{equation*}
对应 $(\calE,\varphi)\mapsto(a_{i}(\varphi))_{1\leq i\leq r}$ 定义了 $\calM_{G,\calL}$ 的{\em Hitchin 纤维化}
\begin{equation*}
f=f_{G,\calL}: \calM_{G,\calL}\to \calA_{G,\calL}:=\prod_{i=1}^{r}\cohog{0}{X,\calL^{\otimes d_{i}}}.
\end{equation*}
目标 $\calA_{G,\calL}$ 同样看作 $k$ 上仿射空间，称为{\em Hitchin 基}。

定义 Hitchin 基 $\calA_{G,\calL}$ 更内蕴的方式如下。仿射概形 $\frc=\Spec \Sym(\frg^{*})^{G}$ 带有 $\Gm$-作用，它诱导坐标环上的分次。设 $\Tot^{\times}(\calL)\to X$ 是 $\calL$ 全空间中零截面的补集。考虑 $\frc$ 在 $X$ 上的 $\calL$-扭版本：
\begin{equation*}
\frc_{X,\calL}:=(\frc\times \Tot^{\times}(\calL))/\Gm,
\end{equation*}
其中 $\l\in\Gm$ 按 $\l: (c,\wt{x})\mapsto (\l c, \l^{-1}\wt{x})$ 作用在两个坐标上。这是 $X$ 上的仿射空间丛，纤维同构于 $\frc$。于是 $\calA_{G,\calL}$ 可典范地等同于映射 $\frc_{X,\calL}\to X$ 的截面构成的模空间。特别地，每个点 $a\in\calA_{G,\calL}$ 给出映射 $a:X\to [\frc/\Gm]$。

\subsubsection{一般正则半单的轨迹} 把 $\calL$ 在 $X$ 的一般点 $\eta$ 处平凡化并限制 $a_{i}$ 到 $\eta$，得多项式 $P_{a}(y)=y^{n}-a_{1}y^{n-1}+a_{2}y^{n-2}+\cdots +(-1)^{n}a_{n}\in F[y]$，其中 $F=k(X)$ 是 $X$ 的函数域。当 $P_{a}(y)$ 是 $F[y]$ 中的可分多项式时，称这样的 $a$ {\em 一般正则半单}。$a$ 的一般正则半单性等价于判别式 $\Delta(P_{a})\in\cohog{0}{X,\calL^{n(n-1)}}$ 不消没，因此它定义了开子概形 $\Ah_{n,\calL}\subset \calA_{n,\calL}$。

对一般 $G$，把 $a\in\calA_{G,\calL}$ 看作映射 $a:X\to [\frc/\Gm]$（见 \S\ref{sss:Hit fib general} 末）。称 $a$ {\em 一般正则半单}，若 $a$ 把 $X$ 的一般点映入开子叠 $[\frc^{\rs}/\Gm]$。这定义了开子概形 $\Ah_{G,\calL}\subset\calA_{G,\calL}$，推广了上面对 $\Ah_{n,\calL}$ 的构造。

\subsubsection{几何性质}\label{sss:geom Hit} 当 $\deg \calL>2g-2$ 时，叠 $\calM|_{\Ah_{G,\calL}}$ 光滑，见 \cite[定理 4.14.1]{NgoFL}。此时态射 $f_{G,\calL}$ 在 $\Ah_{G,\calL}$ 上平坦，见 \cite[推论 4.16.4]{NgoFL}。当 $G$ 半单时，还有更开稠的子集 $\calA^{\textup{ani}}_{G,\calL}\subset \Ah_{G,\calL}$，在其上 $\calM$ 是 Deligne--Mumford 叠且映射 $f_{G,\calL}$ 真，见 \cite[命题 6.1.3]{NgoFL}。与仿射 Springer 纤维的几何所涉及的无穷维相比，Hitchin 纤维化的几何性质好得多，然而它与仿射 Springer 纤维密切相关，这一点将在 \S\ref{ss:rel Hit ASF} 看到。

\subsubsection{推广}\label{sss:general Hitchin} 设 $H$ 是 $k$ 上约化群，处于约化群的短正合列
\begin{equation*}
1\to H_{1}\to H\to A\to 1.
\end{equation*}
设 $(V,\rho)$ 是 $H$ 的表示。则可考虑数对 $(\calE,\psi)$，其中 $\calE$ 是 $X$ 上 $H$-主齐性空间，$\psi$ 是相伴丛 $\rho(\calE)$ 的截面。等价地，这样的数对等同于态射 $X\to [V/H]$。可以证明存在代数叠 $\calM_{H, \rho}$ 分类这样的数对。每个 $H$-主齐性空间诱导 $A$-主齐性空间，故有态射 $\alpha: \calM_{H, \rho}\to \Bun_{A}$。固定 $X$ 上 $A$-主齐性空间 $\calE_{A}$，把原像 $\alpha^{-1}(\calE_{A})$ 记为 $\calM_{H,\rho,\calE_{A}}$。

为回到 Hitchin 模叠，考虑 $H=G\times\Gm$、$H_{1}=G$、$A=\Gm$ 的情形。取 $V=\frg$，$H$ 的作用 $\rho$ 定义如下：$G$ 按伴随表示作用，$\Gm$ 按伸缩作用在 $V$ 上。$H$-主齐性空间是由 $G$-主齐性空间与 $X$ 上线丛 $\calL$ 构成的数对。固定线丛 $\calL$（这等价于固定 $A=\Gm$-主齐性空间），得到同构
\begin{equation*}
\calM_{H,\rho, \calL}\cong\calM_{G, \calL}.
\end{equation*}

对如上一般的 $(1\to H_{1}\to H\to A\to 1,V,\rho,\calE_{A})$，可如下定义 Hitchin 基的类比。令 $\frc_{H,\rho}=V\sslash H_{1}$。这是 $\frc$ 的类比，它带有 $A$ 的作用。再作 $\frc_{H,\rho}$ 在 $X$ 上的扭版本
\begin{equation*}
\frc_{H,\rho,\calE_{A}}:=\calE_{A}\twtimes{A}\frc_{H,\rho}.
\end{equation*}
然后把 $\calA_{H,\rho,\calE_{A}}$ 定义为映射 $\frc_{H,\rho,\calE_{A}}\to X$ 的截面构成的模空间。态射 $[V/H]\to [\frc_{H,\rho}/A]$ 进而诱导 Hitchin 纤维化的类比
\begin{equation*}
f_{H,\rho,\calE_{A}}: \calM_{H,\rho,\calE_{A}}\to \calA_{H,\rho,\calE_{A}}.
\end{equation*}

\subsubsection{例子}\label{sss:MHilb} 考虑 $H=\GL(U)\times\Gm\times\Gm$ 与 $V=\End(U)\oplus U^{*}$。作用 $\rho(g,s_{1}, s_{2})$ 在 $V$ 上由 $(A,u^{*})\mapsto (s_{1}gAg^{-1}, s_{2}gu^{*})$ 给出。模叠 $\calM_{H,\rho}$ 通过只记住两个 $\Gm$-主齐性空间而映入 $\Pic(X)\times\Pic(X)$。固定 $(\calL_{1},\calL_{2})\in \Pic(X)\times\Pic(X)$，它在 $\calM_{H,\rho}$ 中的原像分类三元组 $(\calU,\varphi,\beta)$，其中 $(\calU,\varphi)$ 是 $X$ 上秩 $n=\dim V$ 的 Higgs 丛，$\beta$ 是 $\calO_{X}$-线性映射 $\calU\to \calL_{2}$。此情形下的 Hitchin 基与经典 Hitchin 基 $\calA_{n,\calL_{1}}$ 相同。\S\ref{sss:Hilb} 将把此模空间与曲线的 Hilbert 概形联系起来。
